% Section: Partial results at twenty-five atoms (t=5) and two certified
% near-misses.
%
% Sources: WALL_ATTACK_R47/R48/R49/R50 archives, LOOP_STATE TICK-109/110,
% PROGRESS ledger 2026-07-12.  Every computer-assisted claim carries its
% verifier and SHA-256 prefix; see the accompanying claims manifest
% t5_partial_claims.md.
%
% Assumes the companion-paper preamble: amsthm environments theorem/
% proposition/lemma/corollary/definition/remark, booktabs, enumitem.

\section{Partial results at twenty-five atoms, and two certified
near-misses}\label{sec:t5partial}

At parameter \(t\) the configurations of this paper live in support
circuits with \(t^2\) atoms on \(t^2-1\) support edges. The cases
\(t=3\) and \(t=4\) (nine and sixteen atoms) are closed earlier in the
paper. This section collects what is known at \(t=5\): twenty-five
atoms on twenty-four support edges. The results are partial. We prove
two structural theorems by hand (an order window and a selected-support
bound), report an exact machine exclusion of support orders fifteen and
sixteen, and present two eighteen-vertex circuits that pass every local
test and are excluded only by explicit certificates. Orders seventeen
through twenty-one remain open, and nothing in this section touches the
Erd\H{o}s \(n^2/25\) conjecture itself.

\subsection{Support circuits, selections, and owners}

Recall the footprint objects. A cut \(V(G)=V_0\mathbin{\dot\cup}V_1\)
of a triangle-free graph has crossing edges \(B\) and monochromatic
edges \(M\); each \(xy\in M\) with \(d_B(x,y)=4\) carries the support
\(\operatorname{supp}_B(xy)\subseteq B\) of its shortest geodesics. An
inclusion-minimal family of bad edges violating Hall's condition
projects, by the footprint lemma of the companion classification
paper, to the following abstract object. All bipartite graphs below
come with a fixed bipartition into two \emph{classes}.

\begin{definition}\label{def:t5circuit}
Let \(t\ge3\). A \emph{support circuit at parameter \(t\)} is a pair
\((F,\mathcal A)\) where \(F\) is a finite connected bipartite graph
with \(|E(F)|=t^2-1\), and \(\mathcal A\) is a set of \(t^2\) unordered
pairs of vertices of \(F\), the \emph{atoms}, such that:
\begin{enumerate}[label=\textup{(\alph*)}]
\item every atom \(\{x,y\}\) joins two vertices of a common class with
\(d_F(x,y)=4\); writing \(R(a)\) for the set of length-four
\(x\)--\(y\) geodesics of \(F\) (the \emph{rows} of \(a\)) and
\(P(a)\) for the union of their edge sets, the sets \(P(a)\) cover
\(E(F)\) and every edge of \(F\) lies in at least two of them;
\item the \emph{atom graph} \((V(F),\mathcal A)\) is triangle-free;
\item every proper subfamily \(\mathcal A'\subsetneq\mathcal A\) admits
distinct representatives: an injection \(a\mapsto e_a\in P(a)\).
\end{enumerate}
\end{definition}

By Hall's theorem, condition (c) together with
\(|E(F)|=|\mathcal A|-1\) says exactly that the supports violate
Hall's condition with defect one while every proper subfamily
satisfies it: support circuits are the inclusion-minimal Hall
violations, at \(m=t^2\) atoms. At \(t=3\) the unique support circuit
is the nine-atom double star of the companion classification
\cite{Ferudun26supports}. Note
that the graphs \(G_t\) of \cite{Ferudun26supports} are themselves
\emph{not} support circuits for
\(t\ge4\): their violation has defect far exceeding one, and the
nine-atom circuit arises among their minimal subsystems
\cite{Ferudun26supports}.

\begin{definition}\label{def:selection}
A \emph{selection} \(\omega\) assigns to each atom \(a\) one row
\(\omega(a)\in R(a)\). The \emph{selected} edge set is
\(S_\omega=\bigcup_{a}E(\omega(a))\), and \(L_\omega=E(F)\setminus
S_\omega\) is the \emph{latent} edge set.
\end{definition}

\begin{definition}\label{def:owner}
Let \((F,\mathcal A)\) be a support circuit at parameter \(t\), let
\(v\) be a vertex with
\(\deg_F(v)=\deg_{\mathcal A}(v)=t\), and let \(x_0\in N_F(v)\). A
selection \(\omega\) \emph{realizes a covered rotating owner at
\((v,x_0)\)} if
\begin{enumerate}[label=\textup{(\roman*)}]
\item the only selected rows containing \(v\) are the rows of the
\(t\) atoms at \(v\);
\item \(vx_0\notin S_\omega\) (the \emph{active} edge);
\item every other edge at \(v\) lies in \(S_\omega\);
\item for every \(y\in N_F(v)\setminus\{x_0\}\), some selected row
contains both \(x_0\) and \(y\).
\end{enumerate}
The pairs \(\{x_0,y\}\), \(y\in N_F(v)\setminus\{x_0\}\), are the
\emph{star pairs}; condition (iv) is the \emph{covered star}.
\end{definition}

\begin{definition}\label{def:live}
Let \(\omega\) realize a covered rotating owner at \((v,x_0)\). The
\emph{tail} is the connected component of \(x_0\) in the latent graph
\((V(F),L_\omega)\) with the vertex \(v\) deleted. The owner is
\emph{live} if the tail contains a vertex \(b\) with
\(\{v,b\}\in\mathcal A\), or both members of some atom; otherwise it
is \emph{dead}.
\end{definition}

\begin{definition}\label{def:rotor}
A \emph{rotor core} in \((F,\mathcal A)\) consists of an atom
\(\{a,b\}\), two further vertices \(v\neq m\) in the class of \(a\),
and vertices \(x,y\) in the opposite class, such that both
\((a,x,v,y,b)\) and \((a,x,m,y,b)\) are rows of \(\{a,b\}\). A
\emph{balanced rotor at parameter \(t\)} is a support circuit at
parameter \(t\) carrying a rotor core whose middle vertices \(v,m\)
both satisfy \(\deg_F=\deg_{\mathcal A}=t\), at least one of which
admits a selection realizing a \emph{live} covered rotating owner.
\end{definition}

The full rotor object of the programme carries additional bookkeeping
on top of Definition~\ref{def:rotor}; every such object contains the
data listed here, so all exclusions below apply to it a fortiori.

We record four elementary lemmas. Throughout,
\((F,\mathcal A)\) is a support circuit at parameter \(t\) and
\(v,x_0\) are as in Definition~\ref{def:owner}.

\begin{lemma}[Tightness]\label{lem:tight}
A selection satisfying \textup{(i)} exists if and only if every atom
not at \(v\) has a row avoiding \(v\). If some atom not at \(v\) has
all of its rows through \(v\), then every selection has at least
\(t+1\) selected rows through \(v\).
\end{lemma}

\begin{proof}
Every row of an atom at \(v\) is a geodesic from \(v\), so it contains
\(v\); this gives \(t\) selected rows through \(v\) under any
selection. If every atom not at \(v\) has a \(v\)-avoiding row, choose
such rows for all of them; (i) holds. Conversely an atom not at \(v\)
all of whose rows contain \(v\) contributes a \((t+1)\)-st selected
row through \(v\) under every selection.
\end{proof}

\begin{lemma}[Positions]\label{lem:positions}
Let \(\omega\) satisfy \textup{(i)} and \textup{(ii)}. Then no
selected row of an atom at \(v\) contains \(x_0\), and every selected
row containing \(x_0\) avoids \(v\).
\end{lemma}

\begin{proof}
A row of an atom at \(v\) is a geodesic \((v,q_1,q_2,q_3,q_4)\), so
\(d_F(v,q_k)=k\). The vertex \(x_0\) lies in the class opposite
\(v\), hence could only be \(q_1\) or \(q_3\). As \(q_1\) it would put
the edge \(vx_0\) into \(S_\omega\), against (ii); as \(q_3\) it would
force \(d_F(v,x_0)=3\), against \(vx_0\in E(F)\). For the second
statement, a selected row through \(v\) is a row of an atom at \(v\)
by (i), and therefore omits \(x_0\).
\end{proof}

\begin{lemma}[Step surjection]\label{lem:steps}
Let \(\omega\) satisfy \textup{(i)}--\textup{(iii)}. Then the selected
row of each atom at \(v\) has \(v\) as an endpoint and its first edge
is \(vy\) for some \(y\neq x_0\); the resulting map from the \(t\)
atoms at \(v\) to the \(t-1\) edges \(vy\) is surjective, taking one
value twice and every other value once.
\end{lemma}

\begin{proof}
A row of an atom at \(v\) contains \(v\) exactly once, as an endpoint,
so its unique edge at \(v\) is its first edge, and (ii) rules out
\(vx_0\). Conversely each edge \(vy\), \(y\neq x_0\), is selected by
(iii); a selected row using \(vy\) contains \(v\), is a row of an atom
at \(v\) by (i), and has \(vy\) as its first edge. So the map is onto,
and \(t\) atoms mapping onto \(t-1\) edges with every fibre nonempty
forces fibre sizes \((2,1,\dots,1)\).
\end{proof}

\begin{lemma}[Parity]\label{lem:parity}
Let \(Q\) be a row of some atom with \(v\notin Q\). Then \(Q\)
contains at most one star pair.
\end{lemma}

\begin{proof}
Two star pairs contained in \(Q\) would give three distinct vertices
of \(N_F(v)\) on \(Q\). All of \(N_F(v)\) lies in the class opposite
\(v\), and the vertices of the path \(Q=(q_0,\dots,q_4)\) alternate
classes, so members of \(N_F(v)\) occupy positions of a single parity;
three of them must occupy positions \(0,2,4\). Then \(q_0,q_4\in
N_F(v)\) are the endpoints of the atom of \(Q\), and
\(q_0\,v\,q_4\) is a path in \(F\), so \(d_F(q_0,q_4)\le2\),
contradicting the atom distance four.
\end{proof}

\begin{corollary}[Coverage; order upper bound]\label{cor:coverage}
Let \(\omega\) realize a covered rotating owner at \((v,x_0)\), with
\(t=5\). Then there are four pairwise distinct atoms
\(a_1,\dots,a_4\), none at \(v\), whose selected rows
\(Q_i=\omega(a_i)\) satisfy \(v\notin Q_i\) and
\(\{x_0,y_i\}\subseteq Q_i\), one for each star pair. On \(Q_i\) the
vertices \(x_0\) and \(y_i\) have a common neighbour \(q_i\), and the
four \(4\)-cycles \(v\,x_0\,q_i\,y_i\) are linearly independent in the
cycle space of \(F\). Consequently the cycle rank of \(F\) is at least
four and
\[
|V(F)|\;\le\;24-4+1\;=\;21 .
\]
\end{corollary}

\begin{proof}
By (iv) each star pair lies on a selected row, which avoids \(v\) by
Lemma~\ref{lem:positions}; a row avoiding \(v\) is not a row of an atom
at \(v\). Each atom contributes one selected row, and by
Lemma~\ref{lem:parity} that row contains at most one star pair, so the
four star pairs are covered by four distinct atoms. On \(Q_i\) the
vertices \(x_0,y_i\) occupy positions of equal parity; positions
\(\{0,4\}\) are impossible (the endpoints of \(Q_i\) form the atom
\(a_i\), and two neighbours of \(v\) are at distance at most two), so
the positions differ by two and there is a common row-neighbour
\(q_i\neq v\). The cycle \(C_i=v\,x_0\,q_i\,y_i\) uses the edges
\(vx_0,\,x_0q_i,\,q_iy_i,\,y_iv\) of \(F\). Each \(C_i\) contains
\(vy_i\), while \(C_j\) for \(j\neq i\) meets \(v\) only in
\(vx_0,vy_j\); so each cycle owns a private edge and the four are
independent. For a connected graph the cycle rank is
\(|E|-|V|+1\).
\end{proof}

\subsection{The order window \texorpdfstring{\(15\le|V(F)|\le21\)}{15..21}}

Fix \(t=5\), a support circuit \((F,\mathcal A)\), and a selection
\(\omega\) realizing a covered rotating owner at \((v,x_0)\). Write
\[
K=\{v\}\cup N_F(v)\cup N_{\mathcal A}(v),\qquad |K|=11,
\]
where \(N_{\mathcal A}(v)\) is the set of five atom-partners of
\(v\) (they lie in the class of \(v\) and are distinct from \(v\);
\(N_F(v)\) lies in the opposite class). Vertices outside \(K\) are
\emph{external}.

\begin{lemma}[Endpoint trichotomy]\label{lem:trichotomy}
Let \(Q\) be a covering row for the star pair \(\{x_0,y\}\), that is,
a row of some atom with \(v\notin Q\supseteq\{x_0,y\}\). Then every
vertex of \(Q\) in the class of \(v\) that is adjacent on \(Q\) to a
member of \(N_F(v)\) is external, and exactly one of the following
holds.
\begin{enumerate}[label=\textup{(\Roman*)}]
\item \(x_0\) and \(y\) are interior on \(Q\); then the atom of \(Q\)
joins two external vertices of the class of \(v\), and \(Q\) contains
three external vertices of that class.
\item One of \(x_0,y\) is an endpoint of \(Q\); then the other
endpoint of \(Q\) is an external vertex of the class opposite \(v\),
and \(Q\) contains two external vertices of the class of \(v\).
\end{enumerate}
\end{lemma}

\begin{proof}
A vertex \(w\) of \(Q\) in the class of \(v\), adjacent on \(Q\) to
some \(z\in N_F(v)\), satisfies \(d_F(v,w)\le2\); it is not \(v\)
because \(v\notin Q\), and not in \(N_{\mathcal A}(v)\) because atoms
have distance four. So \(w\) is external. As in
Lemma~\ref{lem:parity}, \(x_0\) and \(y\) occupy positions of equal
parity on \(Q=(q_0,\dots,q_4)\), and not \(\{0,4\}\); so the positions
are \(\{1,3\}\) (case I) or \(\{0,2\}\), \(\{2,4\}\) (case II).

Case I: positions \(0,2,4\) lie in the class of \(v\), and each of
\(q_0,q_2,q_4\) is adjacent on \(Q\) to \(q_1\) or \(q_3\), i.e.\ to a
member of \(N_F(v)\); all three are external, and the atom of \(Q\) is
\(\{q_0,q_4\}\).

Case II, say positions \(\{0,2\}\): now \(q_1,q_3\) lie in the class
of \(v\) and each is adjacent on \(Q\) to \(q_0\) or \(q_2\), members
of \(N_F(v)\); both are external. The atom of \(Q\) is
\(\{q_0,q_4\}\) with \(q_0\in\{x_0,y\}\subseteq N_F(v)\). If \(q_4\in
N_F(v)\) then \(d_F(q_0,q_4)\le2\), against the atom distance; and
\(q_4\) lies in the class opposite \(v\), whose part of \(K\) is
exactly \(N_F(v)\). So \(q_4\) is external.
\end{proof}

\begin{theorem}[Order lower bound]\label{thm:order15}
Let \((F,\mathcal A)\) be a support circuit at \(t=5\) admitting a
selection that realizes a covered rotating owner. Then
\(|V(F)|\ge15\). Together with Corollary~\ref{cor:coverage},
\[
15\;\le\;|V(F)|\;\le\;21 .
\]
\end{theorem}

\begin{proof}
By Corollary~\ref{cor:coverage} there are four covering rows belonging
to four distinct atoms. By Lemma~\ref{lem:trichotomy} each covering
row contributes at least two external vertices in the class of \(v\);
a row of type I contributes three, and if all four rows are of type
II, their atoms have external endpoints in the opposite class, so at
least one external vertex lies there. In either case there are at
least three external vertices and \(|V(F)|\ge11+3=14\).

Suppose \(|V(F)|=14\), so there are exactly three external vertices.

\emph{Case 1: some covering row has type I.} Its three external
vertices lie in the class of \(v\) and exhaust the external supply, so
no external vertex lies in the opposite class, and by
Lemma~\ref{lem:trichotomy} no covering row has type II. Hence all four
covering atoms have type I: four \emph{distinct} atoms, each a pair of
external vertices of the class of \(v\). But only three such vertices
exist, admitting at most \(\binom32=3\) pairs. Contradiction.

\emph{Case 2: all four covering rows have type II.} Then at least two
external vertices lie in the class of \(v\) and at least one in the
opposite class; with three externals in total this forces exactly
\(2+1\), and the classes of \(F\) have sizes
\[
\underbrace{1+5+2}_{\text{class of }v}=8,
\qquad
\underbrace{5+1}_{\text{opposite class}}=6 .
\]
Atoms join vertices of a common class and the atom graph is
triangle-free, so by Mantel's theorem
\[
25=|\mathcal A|\;\le\;\Bigl\lfloor\tfrac{8^2}{4}\Bigr\rfloor
+\Bigl\lfloor\tfrac{6^2}{4}\Bigr\rfloor=16+9=25 .
\]
Equality forces the atom graph induced on the class of \(v\) to be the
unique extremal graph \(K_{4,4}\), in which every vertex has
atom-degree four. But \(\deg_{\mathcal A}(v)=5\). Contradiction.
\end{proof}

\subsection{The selected-support bound}

\begin{theorem}\label{thm:selected}
Let \(t\ge3\), let \((F,\mathcal A)\) be a support circuit at
parameter \(t\), and let \(\omega\) realize a covered rotating owner
at \((v,x_0)\). Then
\[
|S_\omega|\;\ge\;3t-1,
\qquad\text{hence}\qquad
|L_\omega|\;\le\;(t^2-1)-(3t-1)\;=\;t(t-3).
\]
\end{theorem}

\begin{proof}
Write \(Y=N_F(v)\setminus\{x_0\}\), \(|Y|=t-1\), and let
\(b_1,\dots,b_t\) be the atom-partners of \(v\). We exhibit \(3t-1\)
distinct selected edges.

\emph{Star edges.} The \(t-1\) edges \(vy\), \(y\in Y\), are selected
by (iii).

\emph{Second edges.} For each \(y\in Y\) pick, by
Lemma~\ref{lem:steps}, a selected row \(\rho^y\) of an atom at \(v\)
with first edge \(vy\), and let \(\sigma(y)\) be its second edge,
joining \(y\) to a vertex of the class of \(v\). Distinct \(y\) give
distinct edges \(\sigma(y)\): they have distinct endpoints in the
class opposite \(v\). This yields \(t-1\) selected edges not incident
with \(v\) (a row of an atom at \(v\) meets \(v\) only in its first
edge).

\emph{Last edges.} For each \(j\) let \(\tau(b_j)\) be the last edge
of the selected row of the atom \(\{v,b_j\}\), joining \(b_j\) to a
vertex \(q_j\) of the opposite class. The \(b_j\) are pairwise
distinct, so these are \(t\) distinct selected edges, again not
incident with \(v\).

\emph{Distinctness.} The star edges meet \(v\); the others do not. If
\(\sigma(y)=\tau(b_j)\) for some \(y,j\), then matching endpoints by
class gives \(q_j=y\in N_F(v)\) and \(b_j\) adjacent to \(y\), whence
\(d_F(v,b_j)\le2\), contradicting the atom distance four. So the three
families are pairwise disjoint: \(3t-2\) selected edges so far. None
of them contains \(x_0\): the star edges have endpoints
\(v\) and \(y\neq x_0\); each \(\sigma(y)\) joins \(y\neq x_0\) to the
class of \(v\); each \(\tau(b_j)\) joins \(b_j\) to \(q_j\), and
\(q_j\neq x_0\) because rows of atoms at \(v\) avoid \(x_0\)
(Lemma~\ref{lem:positions}).

\emph{One more edge at \(x_0\).} By (iv) some selected row \(P\)
contains \(x_0\) (any covered star pair supplies one); by
Lemma~\ref{lem:positions} \(v\notin P\), so \(P\) has an edge at
\(x_0\) different from \(vx_0\). It is selected, contains \(x_0\), and
therefore differs from all \(3t-2\) edges above. Hence
\(|S_\omega|\ge3t-1\), and the latent bound follows from
\(|E(F)|=t^2-1\).
\end{proof}

\begin{corollary}\label{cor:latent}
For \(t=5\): at least \(14\) of the \(24\) support edges are selected,
at most \(10\) are latent, and, since \(vx_0\in L_\omega\) while every
other edge at \(v\) is selected, the tail of
Definition~\ref{def:live} spans at most \(9\) latent edges.
\end{corollary}

\begin{remark}\label{rem:t3}
At \(t=3\) Theorem~\ref{thm:selected} gives \(|L_\omega|\le0\), while
(ii) forces \(vx_0\in L_\omega\). So no covered rotating owner exists
at nine atoms at all --- consistent with the unique nine-atom circuit,
in which no vertex has support- and atom-degree both equal to three.
At \(t=4\) the latent budget is \(4\); at \(t=5\) it is \(10\). The
budget \(t(t-3)\) grows quadratically, so the tail mechanism exploited
below is specific to small \(t\).
\end{remark}

\begin{remark}
Lemmas~\ref{lem:tight}--\ref{lem:parity} convert the existence of a
covered rotating owner at \((v,x_0)\) into a finite test on the row
database of the circuit: every atom off \(v\) needs a \(v\)-avoiding
row; every atom at \(v\) needs a row whose first edge avoids
\(x_0\); the atom-to-first-edge relation needs a matching covering the
\(t-1\) non-active star edges; and, by Lemma~\ref{lem:parity}, the
covered star needs \(t-1\) \emph{distinct} atoms, one covering row
each, so the pair-to-atom coverage relation needs a matching
saturating the star pairs. Conversely, choosing rows along such
matchings (arbitrary \(v\)-avoiding rows elsewhere) realizes the
owner, so the four conditions are exactly equivalent to realizability.
This four-part test is what the enumeration below implements.
\end{remark}

\subsection{Exact exclusion of support orders fifteen and sixteen}

The machine layer works with a relaxation. A balanced rotor
(Definition~\ref{def:rotor}) whose support has \(n\) vertices, with
\(p\) of them in the class of the owners, satisfies the following
finite constraint system \((R_{p,n-p})\):
\begin{enumerate}[label=\textup{(R\arabic*)}]
\item \(F\) is connected and bipartite with classes of sizes
\(p\) and \(n-p\) and \(|E(F)|=24\);
\item there are vertices \(v,m,a,b\) in the first class and \(x,y\) in
the second with \((a,x,v,y,b)\) and \((a,x,m,y,b)\) paths in \(F\);
\item \(\deg_F(v)=\deg_F(m)=5\);
\item every vertex of \(F\) has positive degree;
\item \(d_F(a,b)=4\);
\item \(v\) and \(m\) each have at least five vertices of their own
class at distance exactly four;
\item at least \(25\) same-class pairs of \(V(F)\) are at distance
exactly four.
\end{enumerate}
Indeed (R2), (R3) restate the rotor core and owner degrees, (R5) holds
because \(\{a,b\}\) is an atom, (R6) because
\(\deg_{\mathcal A}(v)=\deg_{\mathcal A}(m)=5\) and atoms are
same-class distance-four pairs, and (R7) because the \(25\) atoms are
distinct such pairs.

\begin{lemma}\label{lem:shores}
In a balanced rotor at \(t=5\), the class of the owners has at least
seven vertices and the opposite class at least five.
\end{lemma}

\begin{proof}
The owner class contains \(v\), \(m\), and the five atom-partners of
\(v\); these are seven distinct vertices, since \(d_F(v,m)=2\) (via
\(x\)) while atom-partners of \(v\) are at distance four. The opposite
class contains \(N_F(v)\), of size five.
\end{proof}

So at order \(n\) the owner-class size \(p\) ranges over
\(7\le p\le n-5\): four splits at \(n=15\), five at \(n=16\), six at
\(n=17\).

\begin{proposition}[Computer-assisted]\label{prop:orders1516}
For \(n=15\) and every split
\((p,n-p)\in\{(7,8),(8,7),(9,6),(10,5)\}\), and for \(n=16\) and every
split \((p,n-p)\in\{(7,9),(8,8),(9,7),(10,6),(11,5)\}\), the
constraint system \((R_{p,n-p})\) is infeasible. Consequently no
balanced rotor at \(t=5\) has a support of order \(15\) or \(16\), and
with Theorem~\ref{thm:order15} the support order of a balanced rotor
at \(t=5\) lies in \(\{17,\dots,21\}\).
\end{proposition}

The nine infeasibility certificates were produced by the Boolean
constraint model \texttt{rooted\_t5\_support\_cp\_sat.py} (OR-Tools
CP-SAT; integer and Boolean arithmetic only; the symmetry breaking
sorts the degrees of the non-distinguished vertices, which any
solution can be relabelled to satisfy). Each verdict is a completed
infeasibility proof, not a time-out. Artifact SHA-256 prefixes:
\texttt{14ac9d93}, \texttt{8d194613}, \texttt{66e0813c},
\texttt{5721cbc5} (order \(15\), splits as listed) and
\texttt{8dff49ea}, \texttt{bd0c0320}, \texttt{cc5462d6},
\texttt{b5037e11}, \texttt{9f9401d0} (order \(16\)). Unlike the
certificates of the next subsection, these nine rest on a single
solver and have not been independently replayed; we flag this
honestly.

\subsection{Two certified near-misses at order eighteen}

The enumeration at order eighteen produced two support graphs whose
circuits pass \emph{every} local test of this section. We print them
in graph6 form; vertices \(0\)--\(8\) form the owner class and
\(9\)--\(17\) the opposite class, with
\(v=0\), \(m=1\), \(a=2\), \(b=3\), \(x=9\) in the labelling of
\textup{(R2)}:
\[
\begin{array}{ll}
\text{circuit } \#298: &
\texttt{Q??????wE\_[?EGs?D\_@A?C\_B???}\\[2pt]
\text{circuit } \#264: &
\texttt{Q??????wE\_Bws?s?DCD??@?@???}
\end{array}
\]
The following facts are immediate from the printed strings: each graph
has \(18\) vertices, \(24\) edges, classes of sizes \(9+9\), is
connected and bipartite, satisfies
\(\deg_F(0)=\deg_F(1)=5\) and \(d_F(2,3)=4\), and carries the rotor
core of \textup{(R2)}. The same-class distance-four supply is \(32\)
pairs in \(\#298\) and \(30\) pairs in \(\#264\), from which the
twenty-five atoms are drawn.

\begin{proposition}[Computer-assisted]\label{prop:hits}
For each of the two graphs above there is a set \(\mathcal A\) of
twenty-five of its same-class distance-four pairs and a selection
\(\omega\) such that \((F,\mathcal A)\) is a support circuit at
\(t=5\) --- in particular the atom graph is triangle-free --- and
\(\omega\) realizes a covered rotating owner at \((v,x_0)\), with
\(v=0\) and \(x_0=17\) for \(\#298\), and \(v=0\) and \(x_0=x=9\) for
\(\#264\). In \(\#264\) the active edge is the core edge \(vx\)
itself.
\end{proposition}

The primary search and an independent exact replay
(\texttt{verify\_t5\_local\_classifier\_hit.py}, exact graph
operations and integer bipartite matching, no floating point) agree.
Artifact prefixes: \texttt{c1d474d7} (\(\#298\) hit) with verification
artifact \texttt{48ce1638}, and \texttt{6595501f} (\(\#264\) hit) with
verification artifact \texttt{d9e73413}; the verifier script itself
has SHA-256 prefix \texttt{017d1e44}.
Proposition~\ref{prop:hits} is a negative structural datum: it shows
that no purely local argument --- circuit axioms, covered rotating
owner, triangle-freeness of the atom graph --- can close twenty-five
atoms, since all of these conditions are simultaneously realizable.
Both circuits are nevertheless excluded, by the two certificate layers
that follow. Throughout, \emph{the} circuit \(\#298\) (respectively
\(\#264\)) refers to the exhibited witness of
Proposition~\ref{prop:hits}: the atom set and selection of the pinned
rebuild shipped in the ancillary files
(\texttt{anc/t5\_artifacts/}), over which the mechanical certificates
and the ambient exclusions below quantify. The by-hand proof of
Proposition~\ref{prop:dead} is independent of this choice: it applies
verbatim to every twenty-five-atom circuit on the printed graphs at
the recorded pairs.

\begin{proposition}[Dead owners]\label{prop:dead}
In both circuits, every selection realizing the covered rotating owner
at the recorded pair \((v,x_0)\) leaves the owner dead.
\end{proposition}

\begin{proof}
\emph{Circuit \(\#298\), \((v,x_0)=(0,17)\).} From the printed string,
\(\deg_F(17)=2\) with \(N_F(17)=\{0,1\}\). Let \(\omega\) realize the
owner. Each covering row for a star pair \(\{17,y\}\) contains \(17\)
and avoids \(0\) (Lemma~\ref{lem:positions}), so its edges at \(17\)
lie in \(\{\{1,17\}\}\); hence the edge \(\{1,17\}\) is selected. In
the latent graph minus \(v\), the vertex \(17\) is therefore isolated:
the tail is \(\{17\}\). Since \(17\) lies in the class opposite
\(v\), it is not an atom-partner of \(v\), and no atom lies inside a
single vertex. The owner is dead. Moreover
\(N_F(0)=N_F(1)=\{9,10,11,13,17\}\), so the transposition of \(0\) and
\(1\) is an automorphism of \(F\) and the same argument disposes of
the mirror profile at \((1,17)\).

\emph{Circuit \(\#264\), \((v,x_0)=(0,9)\).} From the printed string,
\(\deg_F(9)=3\) with \(N_F(9)=\{0,1,2\}\), and
\(N_F(0)=\{9,10,12,13,15\}\), so the star pairs are \(\{9,y\}\) with
\(y\in\{10,12,13,15\}\). Exhaustive inspection of the full geodesic
database of \(F\) (a finite check on the printed graph) shows:
\begin{enumerate}[label=\textup{(f\arabic*)}]
\item every row of \emph{any} same-class distance-four pair that
contains \(9\) and one of \(10,12,13,15\) and avoids \(0\) uses the
edge \(\{1,9\}\) (the pairs \(\{9,10\},\{9,12\},\{9,13\}\) each have
exactly three such rows, \(\{9,15\}\) exactly one);
\item the unique such row through \(9\) and \(15\) is
\((15,2,9,1,17)\), a row of the pair \(\{15,17\}\), and it uses the
edge \(\{2,9\}\) as well.
\end{enumerate}
Let \(\omega\) realize the owner. Covering rows avoid \(0\)
(Lemma~\ref{lem:positions}), so by (f1) the edge \(\{1,9\}\) is
selected, and by (f2) the covering row for \(\{9,15\}\) must be
\((15,2,9,1,17)\) --- in particular \(\{15,17\}\in\mathcal A\) --- so
\(\{2,9\}\) is selected too. Thus both edges at \(9\) other than the
active edge \(\{0,9\}\) are selected, the tail is \(\{9\}\), and the
owner is dead as before.
\end{proof}

The database facts (f1), (f2) are also certified mechanically:
per-edge unsatisfiability of the alternative
(\texttt{verify\_t5\_tail\_blanket\_unsat.py}, artifact
\texttt{3720a8c2}). Independently of Proposition~\ref{prop:dead}, the
full statement ``over all admissible row selections, no atom has both
endpoints in the latent component of the owner'' was checked as a
monolithic constraint system twice per circuit, by CP-SAT infeasibility
and by a CaDiCaL replay of a CNF encoding (for \(\#298\): \(1680\)
variables, \(5239\) clauses); artifact prefixes \texttt{f5c0cbca} and
\texttt{a8a160d5} (\(\#298\)), \texttt{79471ef0} and \texttt{c7fbcc70}
(\(\#264\)). These
certificates are \emph{intrinsic}: they quantify over selections of
rows of the circuit itself. Graph-level realizations, whose ambient
crossing graph could a priori supply further structure, are addressed
by the second layer.

\begin{proposition}[Computer-assisted; ambient exclusion of
\(\#264\)]\label{prop:extension}
There is no triangle-free graph \(H\) on at most \(25\) vertices with
a two-colouring of \(V(H)\) extending the classes of circuit
\(\#264\) such that: the crossing edges of \(H\) contain \(E(F)\); the
monochromatic edges of \(H\) are exactly the \(25\) atoms; for every
atom \(\{x,y\}\) the distance of \(x,y\) in the crossing graph of
\(H\) is four and every length-four crossing geodesic is a row of the
circuit; and the two-colouring is a maximum cut of \(H\). The eight
constraint systems (one for each assignment of \(k\in\{0,\dots,7\}\)
added vertices to the owner class, unused vertices being permitted to
be isolated) are all infeasible.
\end{proposition}

The primary run (\texttt{extend\_t5\_hit\_maxcut.py}) separates the
maximum-cut switch inequalities exactly and lazily; an independent SAT
replay (\texttt{verify\_t5\_maxcut\_extension\_unsat.py}) reproduces
all eight infeasibilities from the unrestricted crossing-edge domain.
The obstruction is finite and explicit. The vertex set
\(S=\{4,5,6,7,8,11,14,16\}\) of the circuit is crossed by \(23\)
atoms and only \(2\) support edges, so switching \(S\) forces any
row-preserving maximum-cut extension to add at least \(21\) crossing
edges on the boundary of \(S\); exact counting of addable,
row-preserving edges gives capacities
\(21,19,17,15,13,11,9,7\) across the eight assignments, killing seven
of them outright, and the single assignment with capacity \(21\) is
killed jointly by a second switch (joint demand \(42\) against joint
capacity \(28\)).

\begin{proposition}[Computer-assisted; ambient exclusion of
\(\#298\)]\label{prop:298extension}
There is no triangle-free graph \(H\) on at most \(25\) vertices with
a two-colouring of \(V(H)\) extending the classes of circuit
\(\#298\) such that: the crossing edges of \(H\) contain \(E(F)\); the
monochromatic edges of \(H\) are exactly the \(25\) atoms; for every
atom \(\{x,y\}\) the distance of \(x,y\) in the crossing graph of
\(H\) is four and every length-four crossing geodesic is a row of the
circuit; and the two-colouring is a maximum cut of \(H\). The eight
constraint systems (one for each assignment of \(k\in\{0,\dots,7\}\)
added vertices to the owner class, unused vertices being permitted to
be isolated) are all infeasible.
\end{proposition}

The same two verifiers as for \(\#264\) agree on all eight
infeasibilities. Here the displayed two-colouring already fails to be
a maximum cut of the bare graph \((V(F),E(F)\cup\mathcal A)\): the
vertex set \(S=\{4,5,6,7,8,12,14,15,16\}\) of the circuit is crossed
by \(24\) of the \(25\) atoms and only \(4\) support edges, so
switching \(S\) forces any row-preserving maximum-cut extension to add
at least \(20\) crossing edges on the boundary of \(S\). Unlike for
\(\#264\), this demand alone kills no assignment; the infeasibilities
are joint with the row-preservation clauses and, at
\(k\in\{0,1\}\), a second switch. The certificate chain for
\(\#298\) is included in the ancillary files
(\texttt{anc/t5\_artifacts/298\_extension/}).

\subsection{Status of orders seventeen through twenty-one}

\begin{proposition}[Computer-assisted]\label{prop:order17}
At order \(17\), the systems \((R_{7,10})\), \((R_{8,9})\) and
\((R_{12,5})\) are infeasible (artifact prefixes \texttt{4d450ca0},
\texttt{f76b9d64}, \texttt{325012ec}, in the split order listed, the
mapping being fixed by the bit-exact regeneration in the ancillary
files; single-solver certificates, as
in Proposition~\ref{prop:orders1516}). The systems \((R_{9,8})\),
\((R_{10,7})\), \((R_{11,6})\) are feasible.
\end{proposition}

Feasibility of a support relaxation is not existence of a rotor: at
each of the three feasible splits a staged search of \(3{,}000\)
support graphs produced no circuit realizing a live covered rotating
owner (artifacts \texttt{95dbc901}, \texttt{d612c59e},
\texttt{31d82c62}), and at order \(18\) the first \(350\) supports at
split \((9,9)\) were likewise all rejected at the live-owner gate.
These are bounded searches and prove nothing beyond the instances
examined. Orders \(19\)--\(21\) have not been examined. The honest
summary is:

\begin{itemize}[leftmargin=2em]
\item orders \(\le14\): impossible for any covered rotating owner
(Theorem~\ref{thm:order15}), by hand;
\item orders \(15\), \(16\): impossible for balanced rotors
(Proposition~\ref{prop:orders1516}), computer-assisted;
\item order \(17\): open; owner-class size restricted to
\(\{9,10,11\}\) (Proposition~\ref{prop:order17});
\item order \(18\): open; the only circuits found so far that pass all
local tests are \(\#298\) and \(\#264\), and both are excluded by
certificates (Propositions~\ref{prop:dead}, \ref{prop:extension}
and~\ref{prop:298extension});
\item orders \(19\)--\(21\): open and unexamined; order \(22\) and
beyond: impossible (Corollary~\ref{cor:coverage}).
\end{itemize}

All evidence is consistent with the following statement, which we can
neither prove nor refute.

\begin{definition}
Say that a support circuit at \(t=5\) is \emph{scope-vacuous} if every
selection realizing a covered rotating owner, at every pair
\((v,x_0)\), leaves the owner dead.
\end{definition}

\noindent\textbf{Open question.}
\emph{Is every support circuit at twenty-five atoms scope-vacuous?
In particular, is there no balanced rotor at \(t=5\)?}

\medskip

A positive answer would close the parameter \(t=5\) of the programme
in the same sense as the finite closures at \(t=3\) and \(t=4\); by
Remark~\ref{rem:t3} it would still be one more finite base case, not a
statement for all \(t\). A negative answer --- a live covered rotating
owner on a triangle-free \(25\)-atom circuit of order \(17\)--\(21\),
surviving an ambient extension test --- would refute the rotor route
in its present form. Neither outcome, we stress, resolves the Erd\H{o}s inequality
\(\operatorname{bip}(G)\le N^2/25\): the results of this section
constrain one local proof strategy and produce no new bound on
\(\operatorname{bip}(G)\).

\begin{table}[ht]
\centering
\small
\begin{tabular}{lll}
\toprule
fact & verifier(s) & artifact SHA-256 prefixes\\
\midrule
order 15, 4 splits & \texttt{rooted\_t5\_support\_cp\_sat.py} &
\texttt{14ac9d93} \texttt{8d194613} \texttt{66e0813c}
\texttt{5721cbc5}\\
order 16, 5 splits & \texttt{rooted\_t5\_support\_cp\_sat.py} &
\texttt{8dff49ea} \texttt{bd0c0320} \texttt{cc5462d6}
\texttt{b5037e11} \texttt{9f9401d0}\\
order 17, 3 splits & \texttt{rooted\_t5\_support\_cp\_sat.py} &
\texttt{4d450ca0} \texttt{f76b9d64} \texttt{325012ec}\\
\(\#298\) circuit + owner & primary +
\texttt{verify\_t5\_local\_classifier\_hit.py} &
\texttt{c1d474d7}, verification \texttt{48ce1638}\\
\(\#264\) circuit + owner & primary +
\texttt{verify\_t5\_local\_classifier\_hit.py} &
\texttt{6595501f}, verification \texttt{d9e73413}\\
\(\#298\) all-row dead owner & CP-SAT +
\texttt{verify\_t5\_active\_scope\_unsat.py} &
\texttt{f5c0cbca}, \texttt{a8a160d5}\\
\(\#264\) all-row dead owner & CP-SAT + CaDiCaL replay &
\texttt{79471ef0}, \texttt{c7fbcc70}\\
\(\#264\) facts (f1), (f2) &
\texttt{verify\_t5\_tail\_blanket\_unsat.py} &
\texttt{3720a8c2}\\
\(\#264\) ambient exclusion & \texttt{extend\_t5\_hit\_maxcut.py} +
\texttt{verify\_t5\_maxcut\_extension\_unsat.py} & eight
infeasibilities, replayed\\
\(\#298\) ambient exclusion & \texttt{extend\_t5\_hit\_maxcut.py} +
\texttt{verify\_t5\_maxcut\_extension\_unsat.py} &
source \texttt{ff1c1209}, extension \texttt{2bf166b7}, replay
\texttt{d176ce66}\\
\bottomrule
\end{tabular}
\caption{Computer-assisted certificates of this section. Primary
searches use OR-Tools CP-SAT; independent replays use exact graph
operations (NetworkX) or CaDiCaL via PySAT. All arithmetic is integer
or Boolean. The printed prefixes are prefixes of each artifact's
embedded canonical payload hash. The historical artifacts were session
records of the programme ledger; on 2026-07-17 the set was regenerated
from the archived scripts and is shipped in the ancillary files
(\texttt{anc/t5\_artifacts/}, with a per-artifact record). The twelve
split certificates regenerate bit-exactly to the twelve listed hashes.
The circuit-level artifacts depend on a solver-nondeterministic
enumeration order and do not replay byte-exactly; both circuits were
instead rebuilt on the printed support graphs at the recorded
\((v,x_0)\) pairs (\texttt{rebuild\_t5\_local\_classifier\_hit.py})
and every downstream verdict was reproduced --- in particular the
\(\#298\) scope certificate reproduces at exactly \(1680\) variables
and \(5239\) clauses. The historical prefixes are retained above as
the session record. In an independent adversarial
review pass (2026-07-17), the twelve infeasibility and three
feasibility verdicts of
Propositions~\ref{prop:orders1516} and~\ref{prop:order17} were re-run
from the archived model, the two hits of Proposition~\ref{prop:hits}
were re-derived with an independently written constraint encoding, and
the finite facts of Proposition~\ref{prop:dead} were re-verified by
exhaustive geodesic enumeration.}
\label{tab:t5artifacts}
\end{table}
